\chapter{Riemannian Metrics}
\label{chap:pet-riemannian-metrics}

The conventions and notation in this chapter follow \cite{petersen2016}. Unless
stated otherwise, manifolds are smooth, Hausdorff, second countable, and
connected; disconnected constraint sets are allowed when explicitly indicated.

\section{Riemannian Manifolds and Maps}

\begin{definition}[Riemannian manifold]
  \label{def:pet-ch1-riemannian-manifold}
  \lean{PetersenLib.RiemannianMetric}
  \source{petersen:page-0021}
  A \emph{Riemannian metric} on a manifold $M$ is a family of inner products
  $g_p$ on $T_pM$ such that $p\mapsto g_p(X_p,Y_p)$ is smooth for all smooth
  vector fields $X,Y$. The pair $(M,g)$ is a \emph{Riemannian manifold}. We
  write $g(X,Y)$ pointwise and use $g_M$ when the underlying manifold must be
  explicit. If $\dim M=n$, every $(T_pM,g_p)$ is linearly isometric to
  Euclidean $n$-space.
\end{definition}

The elementary models below fix the metric and map conventions used throughout.

\begin{example}[Euclidean space]
  \label{ex:pet-ch1-euclidean-space}
  \lean{PetersenLib.euclideanMetric}
  \source{petersen:page-0021}
  \uses{def:pet-ch1-riemannian-manifold}
  \leanok
  \dcref{ch1:1.1.1}
  Under $T\mathbb{R}^n\simeq\mathbb{R}^n\times\mathbb{R}^n$, where
  $(p,v)$ represents $\frac{d}{dt}(p+tv)|_{t=0}$, the canonical metric is
  \[
    g_{\mathbb{R}^n}((p,v),(p,w))=v\cdot w.
  \]
\end{example}

\begin{definition}[Riemannian isometry]
  \label{def:pet-ch1-riemannian-isometry}
  \lean{PetersenLib.IsRiemannianIsometry}
  \source{petersen:page-0021}
  \uses{def:pet-ch1-riemannian-manifold}
  \leanok
  A diffeomorphism $F:(M,g_M)\to(N,g_N)$ is a \emph{Riemannian isometry} if
  $F^*g_N=g_M$, equivalently
  \[
    g_N(DF_pv,DF_pw)=g_M(v,w)
  \]
  for every $p\in M$ and $v,w\in T_pM$. Its inverse is again a Riemannian
  isometry.
\end{definition}

\begin{example}[Metrics on inner product spaces]
  \label{ex:pet-ch1-inner-product-space}
  \lean{PetersenLib.innerProductSpaceMetric,PetersenLib.linearIsometryEquiv_isRiemannianIsometry}
  \source{petersen:page-0022}
  \uses{def:pet-ch1-riemannian-manifold, def:pet-ch1-riemannian-isometry, ex:pet-ch1-euclidean-space}
  \leanok
  \dcref{ch1:1.1.2}
  An inner product on a finite-dimensional real vector space $V$ defines the
  constant metric $g((p,v),(p,w))=v\cdot w$. A linear isometry between two
  such spaces of equal dimension is a Riemannian isometry. Consequently every
  $n$-dimensional model of this form is isometric to
  $(\mathbb{R}^n,g_{\mathbb{R}^n})$.
\end{example}

\begin{definition}[Pullback metric and Riemannian immersion]
  \label{def:pet-ch1-pullback-metric}
  \lean{PetersenLib.pullbackMetric,PetersenLib.IsRiemannianImmersion}
  \source{petersen:page-0022}
  \uses{def:pet-ch1-riemannian-manifold}
  \leanok
  Let $F:M\to(N,g_N)$ be an immersion. Its \emph{pullback metric} is
  \[
    (F^*g_N)_p(v,w)=g_N(DF_pv,DF_pw).
  \]
  Injectivity of $DF_p$ makes this positive definite. If a metric $g_M$ is
  already specified, $F$ is a \emph{Riemannian} or \emph{isometric immersion}
  when $g_M=F^*g_N$; the analogous terminology applies to embeddings.
\end{definition}

\begin{example}[Canonical metric on a sphere]
  \label{ex:pet-ch1-sphere}
  \lean{PetersenLib.sphereMetric}
  \source{petersen:page-0022}
  \uses{def:pet-ch1-pullback-metric}
  \leanok
  \dcref{ch1:1.1.3}
  For
  $S^n(R)=\{x\in\mathbb{R}^{n+1}:|x|=R\}$, the inclusion into
  $\mathbb{R}^{n+1}$ induces the canonical metric. The unit sphere
  $S^n=S^n(1)$ with this metric is the standard sphere.
\end{example}

\begin{remark}[Nonlinear isometric immersions]
  \label{rem:pet-ch1-nonlinear-isometric-immersion}
  \lean{PetersenLib.curveIsometricEmbedding}
  \source{petersen:page-0023}
  \uses{def:pet-ch1-pullback-metric, ex:pet-ch1-euclidean-space}
  \leanok
  Every unit-speed curve $c:\mathbb{R}\to\mathbb{R}^2$ is an isometric
  immersion. Thus
  \[
    F(x^1,\ldots,x^k)=(c(x^1),x^2,\ldots,x^k)
      :\mathbb{R}^k\to\mathbb{R}^{k+1}
  \]
  is generally nonlinear and is an isometric immersion, or an embedding when
  $c$ is embedded. For instance $c(t)=(\cos t,\sin t)$ is an immersion, while
  $c(t)=(\log(t+\sqrt{1+t^2}),\sqrt{1+t^2})$ is an embedding. Hence an
  isometric immersion need not preserve ambient endpoint distance.
\end{remark}

When $DF$ is surjective, the metric comparison is made on the horizontal
subspace rather than on the entire tangent space.

\begin{definition}[Riemannian submersion]
  \label{def:pet-ch1-riemannian-submersion}
  \lean{PetersenLib.IsRiemannianSubmersion}
  \source{petersen:page-0023,petersen:page-0024}
  \uses{def:pet-ch1-riemannian-manifold}
  \leanok
  A submersion $F:(M,g_M)\to(N,g_N)$ is a \emph{Riemannian submersion} if
  \[
    DF_p:(\ker DF_p)^\perp\longrightarrow T_{F(p)}N
  \]
  is a linear isometry for every $p\in M$. Equivalently,
  $g_M(v,w)=g_N(DF_pv,DF_pw)$ for horizontal $v,w$, or the adjoint
  $(DF_p)^*:T_{F(p)}N\to T_pM$ preserves inner products.
\end{definition}

\begin{example}[Orthogonal projections]
  \label{ex:pet-ch1-orthogonal-projection}
  \lean{PetersenLib.orthogonalProjectionSubmersion}
  \source{petersen:page-0023}
  \uses{def:pet-ch1-riemannian-submersion}
  \leanok
  \dcref{ch1:1.1.4}
  Every orthogonal projection
  $(\mathbb{R}^n,g_{\mathbb{R}^n})\to(\mathbb{R}^k,g_{\mathbb{R}^k})$, $k<n$,
  is a Riemannian submersion.
\end{example}

\begin{example}[Hopf fibration]
  \label{ex:pet-ch1-hopf-fibration}
  \lean{PetersenLib.hopfMap}
  \source{petersen:page-0023,petersen:page-0024}
  \uses{def:pet-ch1-riemannian-submersion, ex:pet-ch1-sphere}
  \leanok
  \dcref{ch1:1.1.5}
  With $S^3(1)\subset\mathbb{C}^2$ and
  $S^2(1/2)\subset\mathbb{R}\oplus\mathbb{C}$, the map
  \[
    H(z,w)=\Big(\tfrac12(|w|^2-|z|^2),z\bar w\Big)
  \]
  is a Riemannian submersion $S^3(1)\to S^2(1/2)$.
\end{example}
\begin{proof}
  \leanok
  \uses{ex:pet-ch1-hopf-fibration}
  \sketch
  The fibre through $(z,w)$ is
  $\{(e^{i\theta}z,e^{i\theta}w)\}$, with tangent $i(z,w)$. Its orthogonal
  complement is spanned over $\mathbb{R}$ by
  $\lambda(-\bar w,\bar z)$, $\lambda\in\mathbb{C}$. Differentiating the
  extension of $H$ along
  $(z,w)+\lambda(-\bar w,\bar z)$ gives
  \[
    DH_{(z,w)}\big(\lambda(-\bar w,\bar z)\big)
      =\big(2\operatorname{Re}(\bar\lambda zw),
             -\lambda\bar w^2+\bar\lambda z^2\big).
  \]
  This vector and $\lambda(-\bar w,\bar z)$ both have length $|\lambda|$.
  Thus $DH$ is isometric on the horizontal subspace, as required by
  \cref{def:pet-ch1-riemannian-submersion}.
\end{proof}

Positive definiteness can be weakened while retaining nondegeneracy.

\begin{definition}[Pseudo-Riemannian metric and index]
  \label{def:pet-ch1-pseudo-riemannian}
  \lean{PetersenLib.PseudoRiemannianMetric,PetersenLib.pseudoRiemannianIndex}
  \source{petersen:page-0024,petersen:page-0025}
  \uses{def:pet-ch1-riemannian-manifold}
  \leanok
  A \emph{semi-Riemannian} or \emph{pseudo-Riemannian metric} is a smooth
  symmetric bilinear form $g$ on the tangent spaces such that for each nonzero
  $v$ some $w$ satisfies $g(v,w)\ne0$. Each $T_pM$ admits a splitting
  $P\oplus N$ on which $g$ is respectively positive and negative definite.
  Although $P,N$ are not canonical, their dimensions are locally constant;
  on connected $M$, the \emph{index} is the common value $\dim N$.
\end{definition}

\begin{example}[Minkowski space]
  \label{ex:pet-ch1-minkowski}
  \lean{PetersenLib.minkowskiMetric}
  \source{petersen:page-0025}
  \uses{def:pet-ch1-pseudo-riemannian}
  \leanok
  \dcref{ch1:1.1.6}
  For $n=n_1+n_2$ and
  $\mathbb{R}^{n_1,n_2}=\mathbb{R}^{n_1}\times\mathbb{R}^{n_2}$, define
  \[
    g((p,v),(p,w))=v_1\cdot w_1-v_2\cdot w_2.
  \]
  This metric has index $n_2$. The cases $n_1=1$ or $n_2=1$ are the usual
  sign conventions for Minkowski space.
\end{example}

\begin{example}[Hyperbolic space]
  \label{ex:pet-ch1-hyperbolic-space}
  \lean{PetersenLib.hyperbolicSpace}
  \source{petersen:page-0025,petersen:page-0026}
  \uses{def:pet-ch1-pseudo-riemannian, ex:pet-ch1-minkowski, def:pet-ch1-pullback-metric}
  \leanok
  \dcref{ch1:1.1.7}
  Hyperbolic space $H^n(R)$ is the component with $x^{n+1}>0$ of
  \[
    (x^1)^2+\cdots+(x^n)^2-(x^{n+1})^2=-R^2
      \quad\text{in }\mathbb{R}^{n,1},
  \]
  equipped with the induced Minkowski form. This induced form is positive
  definite. We write $H^n=H^n(1)$.
\end{example}
\begin{proof}
  \uses{ex:pet-ch1-hyperbolic-space}
  \leanok
  \sketch
  For $p\in H^n(R)$ and $v\in T_pH^n(R)$, differentiation of the level-set
  equation yields
  \[
    v^{n+1}=\frac{\sum_{i=1}^n v^ip^i}{p^{n+1}}.
  \]
  Hence Cauchy--Schwarz and
  $\sum_i(p^i)^2/(p^{n+1})^2=1-(R/p^{n+1})^2$ imply
  \[
    g(v,v)=\sum_i(v^i)^2-
      \left(\frac{\sum_i v^ip^i}{p^{n+1}}\right)^2
      \ge \left(\frac{R}{p^{n+1}}\right)^2\sum_i(v^i)^2.
  \]
  Equality $g(v,v)=0$ therefore forces all components of $v$ to vanish.
\end{proof}

\section{The Volume Form}

Fixing an orientation turns the metric determinant into an alternating
$n$-linear form.

\begin{definition}[Signed volume]
  \label{def:pet-ch1-signed-volume}
  \lean{PetersenLib.signedVolume}
  \source{petersen:page-0026}
  \uses{def:pet-ch1-riemannian-manifold}
  \leanok
  If $(V,g)$ is oriented and $e_1,\ldots,e_n$ is a positively oriented
  orthonormal basis, define
  \[
    \operatorname{vol}(v_1,\ldots,v_n)
      =\det[g(v_i,e_j)]
      =\det\big([v_1,\ldots,v_n][e_1,\ldots,e_n]^T\big).
  \]
  In the basis $e_i$, this is $\det[v_1,\ldots,v_n]$.
\end{definition}

\begin{proposition}[Independence of the orthonormal basis]
  \label{prop:pet-ch1-signed-volume-basis-independent}
  \lean{PetersenLib.signedVolume_orthonormal_basis_invariant}
  \source{petersen:page-0026}
  \uses{def:pet-ch1-signed-volume}
  \leanok
  The value in \cref{def:pet-ch1-signed-volume} is independent of the chosen
  positively oriented orthonormal basis.
\end{proposition}
\begin{proof}
  \uses{prop:pet-ch1-signed-volume-basis-independent}
  \leanok
  \sketch
  For two such bases $e_i,f_i$, the change-of-basis matrix
  $[f_1,\ldots,f_n][e_1,\ldots,e_n]^T$ is orthogonal with determinant $1$.
  Multiplicativity of the determinant gives
  \[
    \det[g(v_i,f_j)]
      =\det\big([v_1,\ldots,v_n][f_1,\ldots,f_n]^T\big)
      =\det[g(v_i,e_j)].
  \]
\end{proof}

\begin{definition}[Volume form on an oriented manifold]
  \label{def:pet-ch1-volume-form-oriented}
  \lean{PetersenLib.volumeForm}
  \source{petersen:page-0026,petersen:page-0027}
  \uses{prop:pet-ch1-signed-volume-basis-independent}
  \leanok
  On an oriented Riemannian $n$-manifold, the \emph{volume form} is
  \[
    (\operatorname{vol}_g)_p(v_1,\ldots,v_n)=\det[g_p(v_i,e_j)],
  \]
  where $e_1,\ldots,e_n$ is any positively oriented orthonormal basis of
  $T_pM$. It is also denoted $d\,\operatorname{vol}$, without asserting that
  it is exact.
\end{definition}

\begin{definition}[Local volume form]
  \label{def:pet-ch1-volume-form-local}
  \lean{PetersenLib.localVolumeForm}
  \source{petersen:page-0027}
  \uses{def:pet-ch1-volume-form-oriented}
  \leanok
  On a domain $U$ carrying an orthonormal frame $E_1,\ldots,E_n$, declare the
  frame positive and set
  \[
    \operatorname{vol}(X_1,\ldots,X_n)=\det[g(X_i,E_j)].
  \]
  This defines a local volume form even when $M$ is not orientable.
\end{definition}

\begin{remark}[Projection identity for volume]
  \label{rem:pet-ch1-height-base}
  \lean{PetersenLib.volumeForm_projection}
  \source{petersen:page-0027}
  \uses{def:pet-ch1-volume-form-local}
  \leanok
  For an oriented orthonormal frame and any vector $X$,
  \[
    \operatorname{vol}(E_1,\ldots,E_{i-1},X,E_{i+1},\ldots,E_n)
      =g(X,E_i).
  \]
\end{remark}

\begin{remark}[Integration on Riemannian manifolds]
  \label{rem:pet-ch1-integration-riemannian}
  \lean{PetersenLib.integrateOnRiemannianManifold}
  \source{petersen:page-0027}
  \uses{def:pet-ch1-volume-form-local}
  \leanok
  On an oriented Riemannian manifold, $\operatorname{vol}_g$ induces a measure
  $\mu$ and functions are integrated as $\int_M f\,d\mu$. For arbitrary $M$,
  there is an open dense $O\subset M$ with
  $TO\simeq O\times\mathbb{R}^n$; an orthonormal frame orients $O$, and
  integration over $O$ defines integration on $M$ independently of global
  orientability.
\end{remark}

\section{Groups and Riemannian Manifolds}

\subsection{Isometry Groups}

\begin{definition}[Isometry group, isotropy, and homogeneity]
  \label{def:pet-ch1-isometry-group}
  \lean{PetersenLib.IsometryGroup,PetersenLib.IsotropyGroup,PetersenLib.IsHomogeneous}
  \source{petersen:page-0027}
  \uses{def:pet-ch1-riemannian-isometry}
  \leanok
  The isometry group of $(M,g)$ is $\operatorname{Iso}(M,g)$, abbreviated
  $\operatorname{Iso}(M)$ when the metric is understood. The isotropy subgroup
  at $p$ is
  \[
    \operatorname{Iso}_p(M,g)=\{F\in\operatorname{Iso}(M,g):F(p)=p\}.
  \]
  The manifold is \emph{homogeneous} if this group acts transitively on $M$.
\end{definition}

The basic constant-curvature models are homogeneous, with the following full
isometry groups.

\begin{example}[Isometries of Euclidean space]
  \label{ex:pet-ch1-isometry-group-euclidean}
  \lean{PetersenLib.isometryGroup_euclideanSpace}
  \source{petersen:page-0028}
  \uses{def:pet-ch1-isometry-group, def:pet-ch1-riemannian-isometry}
  \leanok
  \dcref{ch1:1.3.1}
  \[
    \operatorname{Iso}(\mathbb{R}^n,g_{\mathrm{eu}})
      =\mathbb{R}^n\rtimes O(n)
      =\{x\mapsto v+Ox:v\in\mathbb{R}^n,\ O\in O(n)\},
  \]
  with unique $v,O$. For every $p$, the isotropy is isomorphic to $O(n)$ and
  $\mathbb{R}^n\simeq\operatorname{Iso}/\operatorname{Iso}_p$. More generally,
  a homogeneous manifold is the quotient of its isometry group by an isotropy
  subgroup.
\end{example}
\begin{proof}
  \leanok
  \uses{ex:pet-ch1-isometry-group-euclidean,rem:pet-ch1-ex-17,rem:pet-ch1-ex-19}
  \sketch
  Affine orthogonal maps are isometries. Conversely, a Riemannian isometry $F$
  preserves curve length by \cref{rem:pet-ch1-ex-17}. Applying
  \cref{rem:pet-ch1-ex-19} to the image of the segment from $x$ to $y$ gives
  $|F(x)-F(y)|\le |x-y|$; applying the same argument to $F^{-1}$ gives
  equality. Mazur--Ulam then makes $F$ affine:
  $F(x)=F(0)+Ox$ with $O\in O(n)$.
\end{proof}

\begin{example}[Isometries of the sphere]
  \label{ex:pet-ch1-isometry-group-sphere}
  \lean{PetersenLib.isometryGroup_sphere}
  \source{petersen:page-0028}
  \uses{def:pet-ch1-isometry-group, ex:pet-ch1-sphere}
  \leanok
  \dcref{ch1:1.3.2}
  \[
    \operatorname{Iso}(S^n(R),g_{S^n(R)})=O(n+1)
      =\operatorname{Iso}_0(\mathbb{R}^{n+1},g_{\mathrm{eu}}).
  \]
  Each isotropy subgroup is isomorphic to $O(n)$, and
  $S^n\simeq O(n+1)/O(n)$.
\end{example}
\begin{proof}
  \uses{ex:pet-ch1-isometry-group-sphere,rem:pet-ch1-ex-20,rem:pet-ch1-ex-17}
  \leanok
  \sketch
  Orthogonal maps restrict to sphere isometries. If $F$ is any sphere
  isometry, it preserves curve length. For $x,y\in S^n(R)$, the minimizing
  great-circle arc in \cref{rem:pet-ch1-ex-20} has length
  $R\arccos(\langle x,y\rangle/R^2)$. Comparing its image under $F$ with the
  minimizing arc between $F(x),F(y)$, and then applying the same comparison to
  $F^{-1}$, yields
  $\langle F(x),F(y)\rangle=\langle x,y\rangle$.
  For an orthonormal basis $e_i$, define $Oe_i=R^{-1}F(Re_i)$. These images are
  orthonormal, so $O\in O(n+1)$, and the preserved pairings imply $F(x)=Ox$
  for every $x\in S^n(R)$.
\end{proof}

\begin{example}[Isometries of hyperbolic space]
  \label{ex:pet-ch1-isometry-group-hyperbolic}
  \lean{PetersenLib.isometryGroup_hyperbolicSpace}
  \source{petersen:page-0028,petersen:page-0029}
  \uses{def:pet-ch1-isometry-group, ex:pet-ch1-hyperbolic-space}
  \leanok
  \dcref{ch1:1.3.3}
  Let
  $O(n,1)=\{L:g(Lv,Lv)=g(v,v)\}$ and let $O^+(n,1)$ be its subgroup preserving
  the upper timelike cone. A bijection of $H^n(R)$ is a Riemannian isometry
  exactly when it is the restriction of an element of $O^+(n,1)$. This group
  acts transitively, the isotropy of $Re_{n+1}$ is $O(n)$, and
  $\operatorname{Iso}(H^n(R))$ is realized by $O^+(n,1)$.
\end{example}
\begin{proof}
  \uses{ex:pet-ch1-isometry-group-hyperbolic,rem:pet-ch1-ex-21,rem:pet-ch1-ex-17}
  \leanok
  \sketch
  An element of $O^+(n,1)$ preserves both $g(x,x)=-R^2$ and the upper sheet,
  and its restriction preserves the induced metric.

  Conversely, let $F\in\operatorname{Iso}(H^n(R))$. The minimizing hyperbola
  of \cref{rem:pet-ch1-ex-21} from $x$ to $y$ has length
  $R\operatorname{arcosh}(-g(x,y)/R^2)$. Length preservation and the inverse
  comparison give $g(F(x),F(y))=g(x,y)$.
  The points
  \[
    W_0=(0,R),\qquad
    W_i=(R\sinh(1)e_i,R\cosh(1)),\quad 1\le i\le n,
  \]
  form a basis of $\mathbb{R}^{n,1}$. Extend $F(W_i)$ linearly to $L$.
  Preservation of pairings on this basis gives $L\in O(n,1)$. For every
  $x\in H^n(R)$, the equalities
  $g(F(x),LW_i)=g(Lx,LW_i)$ for all $i$ imply $F(x)=Lx$ by nondegeneracy.
  Finally every upper timelike vector is a positive multiple of a point of
  $H^n(R)$, so $L$ preserves the upper cone and lies in $O^+(n,1)$.
\end{proof}

\subsection{Lie Groups}

\begin{definition}[Left-invariant metric]
  \label{def:pet-ch1-left-invariant-metric}
  \lean{PetersenLib.leftInvariantMetric}
  \source{petersen:page-0029}
  \uses{def:pet-ch1-riemannian-isometry}
  \leanok
  Left translation identifies $TG\simeq G\times T_eG$. Every inner product on
  $T_eG$ therefore extends uniquely to a metric for which all left translations
  are isometries; these are precisely the \emph{left-invariant metrics} on $G$.
  Different inner products can produce nonisometric metrics.
\end{definition}

\begin{remark}[Metrics on homogeneous quotients]
  \label{rem:pet-ch1-quotient-submersion}
  \lean{PetersenLib.homogeneousQuotientMetric}
  \source{petersen:page-0029}
  \uses{def:pet-ch1-riemannian-submersion, def:pet-ch1-left-invariant-metric}
  \leanok
  Let $H\subset G$ be compact. If right translation by every $h\in H$ is an
  isometry of a metric on $G$, there is a unique metric on $G/H$ for which
  $G\to G/H$ is a Riemannian submersion. If the metric on $G$ is also
  left-invariant, the induced left action of $G$ on $G/H$ is isometric and
  transitive. The formalized quotient-metric result uses the more general
  hypothesis of an isometric action; homogeneity is not separately formalized.
\end{remark}

\begin{example}[Fubini--Study metric]
  \label{ex:pet-ch1-fubini-study}
  \lean{PetersenLib.fubiniStudyMetric,PetersenLib.ComplexProjectiveSpace,PetersenLib.projSphere,PetersenLib.chartCP,PetersenLib.fubiniStudyMetricComplexProjectiveSpace}
  \source{petersen:page-0029}
  \uses{rem:pet-ch1-quotient-submersion, ex:pet-ch1-sphere}
  \leanok
  \dcref{ch1:1.3.4}
  Scalar multiplication by $S^1\subset\mathbb{C}$ acts isometrically on
  $S^{2n+1}(1)\subset\mathbb{C}^{n+1}$. The quotient
  $S^{2n+1}/S^1=\mathbb{CP}^n$ carries the unique \emph{Fubini--Study metric}
  making the quotient map a Riemannian submersion. For $n=1$ this is
  $S^3(1)\to\mathbb{CP}^1=S^2(1/2)$.
\end{example}

\begin{example}[$SU(2)$ and Berger spheres]
  \label{ex:pet-ch1-su2-berger}
  \lean{PetersenLib.suTwoMetric,PetersenLib.bergerSphereMetric}
  \source{petersen:page-0029,petersen:page-0030}
  \uses{def:pet-ch1-left-invariant-metric, ex:pet-ch1-sphere}
  \leanok
  \dcref{ch1:1.3.5}
  There are identifications
  \[
    SU(2)=\left\{\left[\begin{smallmatrix}z&w\\-\bar w&\bar z\end{smallmatrix}\right]
      :|z|^2+|w|^2=1\right\}=S^3(1).
  \]
  The Lie algebra has basis
  \[
    X_1=\operatorname{diag}(i,-i),\quad
    X_2=\left[\begin{smallmatrix}0&1\\-1&0\end{smallmatrix}\right],\quad
    X_3=\left[\begin{smallmatrix}0&i\\i&0\end{smallmatrix}\right].
  \]
  Declaring this left-invariant frame orthonormal gives the round metric on
  $S^3(1)$. Declaring it orthogonal with
  $|X_1|=\varepsilon$ and $|X_2|=|X_3|=1$ gives the Berger metric
  $g_\varepsilon$. Since $X_1$ is tangent to the Hopf circle action,
  $g_\varepsilon$ scales the round metric by $\varepsilon^2$ along Hopf fibres.
\end{example}

\subsection{Covering Maps}

\begin{remark}[Metrics induced by coverings and quotients]
  \label{rem:pet-ch1-covering-metric}
  \lean{PetersenLib.coveringInducedMetric,PetersenLib.quotientMetric}
  \source{petersen:page-0030,petersen:page-0031}
  \uses{def:pet-ch1-riemannian-submersion, def:pet-ch1-pullback-metric}
  \leanok
  A covering $F:M\to N$ is both an immersion and a submersion. Pulling back a
  metric on $N$ makes $F$ a local isometry, and the deck group of a normal
  covering acts isometrically. Conversely, if a group $\Gamma$ acts by
  isometries on $M$ and $N=M/\Gamma$ is a covering quotient, then $N$ has a
  unique metric making $M\to N$ a local isometry.
\end{remark}

\begin{example}[Flat tori]
  \label{ex:pet-ch1-flat-torus}
  \lean{PetersenLib.flatTorus}
  \source{petersen:page-0031}
  \uses{rem:pet-ch1-covering-metric}
  \leanok
  \dcref{ch1:1.3.6}
  The quotient $\mathbb{R}^2/\mathbb{Z}^2$ is the flat square torus
  $S^1\times S^1$ with its product metric, which is bi-invariant. More
  generally, a basis $v_1,v_2$ defines the lattice action
  $x\mapsto x+nv_1+mv_2$ and a flat quotient metric determined up to isometry
  by $|v_1|$, $|v_2|$, and the angle between them.
\end{example}

\begin{example}[Real projective space]
  \label{ex:pet-ch1-real-projective-space}
  \lean{PetersenLib.realProjectiveSpaceCovering}
  \source{petersen:page-0031}
  \uses{rem:pet-ch1-covering-metric, ex:pet-ch1-sphere}
  \leanok
  \dcref{ch1:1.3.7}
  The antipodal isometry of $S^n(1)$ induces the Riemannian covering
  $S^n\to\mathbb{RP}^n$.
\end{example}

\section{Local Representations of Metrics}

\subsection{Einstein Summation Convention}

\begin{convention}[Einstein summation]
  \label{conv:pet-ch1-einstein-summation}
  \lean{PetersenLib.einsteinSummationConvention}
  \source{petersen:page-0031,petersen:page-0032,petersen:page-0033}
  Basis vectors have lower indices and components have upper indices, with a
  repeated upper-lower pair summed:
  $v=v^ie_i$. The dual basis satisfies $e^i(e_j)=\delta^i_j$, so
  $v^i=e^i(v)$. For a linear map $L$, write
  $L(e_j)=L^i_j e_i$; upper indices enumerate rows and lower indices columns.
  In coordinates $x^i$ on a manifold, the coordinate frame and coframe are
  \[
    \partial_i=\frac{\partial}{\partial x^i},\qquad
    dx^i(\partial_j)=\delta^i_j.
  \]
  Thus a curve $c(t)=(x^i(t))$ has intrinsic velocity
  $\dot c=\dot x^i\partial_i$, with $\dot x^i=dx^i(\dot c)$.
\end{convention}

\subsection{Coordinate Representations}

For $1$-forms, juxtaposition denotes the ordered tensor product:
$\theta_1\theta_2(v,w)=\theta_1(v)\theta_2(w)$, which need not equal
$\theta_2\theta_1$.

\begin{definition}[Coordinate representation of a metric]
  \label{def:pet-ch1-coordinate-metric}
  \lean{PetersenLib.metricCoordinateComponents}
  \source{petersen:page-0033}
  \uses{def:pet-ch1-riemannian-manifold, conv:pet-ch1-einstein-summation}
  \leanok
  In coordinates $x=(x^1,\ldots,x^n)$, a Riemannian metric is
  \[
    g=g_{ij}\,dx^i dx^j,
    \qquad g_{ij}=g(\partial_i,\partial_j).
  \]
  The smooth matrix $(g_{ij})$ is symmetric and positive definite. This identity
  follows by writing $v=dx^i(v)\partial_i$ and
  $w=dx^j(w)\partial_j$.
\end{definition}

\begin{example}[Euclidean metric in Cartesian coordinates]
  \label{ex:pet-ch1-euclidean-coords}
  \lean{PetersenLib.euclideanMetricCoordinates}
  \source{petersen:page-0033}
  \uses{def:pet-ch1-coordinate-metric, ex:pet-ch1-euclidean-space}
  \leanok
  \dcref{ch1:1.4.1}
  In the identity chart,
  $g_{\mathbb{R}^n}=\delta_{ij}\,dx^i dx^j=\sum_{i=1}^n(dx^i)^2$.
\end{example}

\begin{example}[Euclidean metric in polar coordinates]
  \label{ex:pet-ch1-polar-coords}
  \lean{PetersenLib.polarCoordinateMetric}
  \source{petersen:page-0034}
  \uses{def:pet-ch1-coordinate-metric, ex:pet-ch1-euclidean-coords}
  \leanok
  \dcref{ch1:1.4.2}
  On the complement of a half-line in $\mathbb{R}^2$, polar coordinates give
  \[
    g=dr^2+r^2d\theta^2.
  \]
  Thus $g_{rr}=1$, $g_{r\theta}=g_{\theta r}=0$, and
  $g_{\theta\theta}=r^2$.
\end{example}
\begin{proof}
  \uses{ex:pet-ch1-polar-coords}
  \leanok
  \sketch
  From $x=r\cos\theta$ and $y=r\sin\theta$,
  \[
    dx=\cos\theta\,dr-r\sin\theta\,d\theta,
    \qquad
    dy=\sin\theta\,dr+r\cos\theta\,d\theta.
  \]
  Expanding $dx^2+dy^2$ cancels the mixed terms and yields
  $dr^2+r^2d\theta^2$.
\end{proof}

\subsection{Frame Representations}

\begin{definition}[Frame representation of a metric]
  \label{def:pet-ch1-frame-metric}
  \lean{PetersenLib.frameRepresentation}
  \source{petersen:page-0034}
  \uses{def:pet-ch1-riemannian-manifold}
  \leanok
  A \emph{frame} on $U\subset M$ is a family of pointwise linearly independent
  vector fields $X_1,\ldots,X_n$. Its dual \emph{coframe}
  $\sigma^1,\ldots,\sigma^n$ satisfies $\sigma^i(X_j)=\delta^i_j$, and
  \[
    g=g_{ij}\sigma^i\sigma^j,
    \qquad g_{ij}=g(X_i,X_j).
  \]
\end{definition}

\begin{example}[Left-invariant coframe representation]
  \label{ex:pet-ch1-left-invariant-coframe}
  \lean{PetersenLib.leftInvariantCoframeMetric}
  \source{petersen:page-0034,petersen:page-0035}
  \uses{def:pet-ch1-frame-metric, def:pet-ch1-left-invariant-metric, ex:pet-ch1-su2-berger}
  \leanok
  \dcref{ch1:1.4.3}
  For a left-invariant frame $X_i$ on a Lie group and its dual coframe
  $\sigma^i$, every left-invariant metric has the form
  $g=g_{ij}\sigma^i\sigma^j$ for a constant positive-definite symmetric matrix
  $(g_{ij})$. If the $X_i$ are orthonormal, then
  $g=\sum_i(\sigma^i)^2$. In particular, the Berger metrics are
  \[
    g_\varepsilon=\varepsilon^2(\sigma^1)^2
      +(\sigma^2)^2+(\sigma^3)^2.
  \]
\end{example}

\begin{example}[Metric on a surface of revolution]
  \label{ex:pet-ch1-surface-of-revolution}
  \lean{PetersenLib.surfaceOfRevolutionMetric}
  \source{petersen:page-0035}
  \uses{def:pet-ch1-frame-metric}
  \leanok
  \dcref{ch1:1.4.4}
  Let $c(t)=(r(t),0,z(t))$, with $r(t)>0$, and revolve it about the $z$-axis:
  \[
    f(t,\theta)=(r(t)\cos\theta,r(t)\sin\theta,z(t)).
  \]
  The induced metric is
  \[
    g=(\dot r^2+\dot z^2)dt^2+r^2d\theta^2.
  \]
  When $c$ has unit speed, this becomes $dt^2+r^2d\theta^2$.
\end{example}
\begin{proof}
  \uses{ex:pet-ch1-surface-of-revolution}
  \leanok
  \sketch
  Substitution into $dx^2+dy^2+dz^2$ gives
  \[
  \begin{aligned}
    dx&=\dot r\cos\theta\,dt-r\sin\theta\,d\theta,\\
    dy&=\dot r\sin\theta\,dt+r\cos\theta\,d\theta,\\
    dz&=\dot z\,dt.
  \end{aligned}
  \]
  The mixed terms cancel, leaving the displayed metric.
\end{proof}

\begin{remark}[Rotationally symmetric metrics]
  \label{rem:pet-ch1-rotationally-symmetric}
  \lean{PetersenLib.rotationallySymmetricMetric}
  \source{petersen:page-0036}
  \uses{ex:pet-ch1-surface-of-revolution}
  \leanok
  A rotationally symmetric metric on $I\times S^1$ is
  $\eta^2(t)dt^2+\rho^2(t)d\theta^2$; a reparametrization of $I$ usually makes
  $\eta=1$. Such a metric need not be induced by a surface of revolution:
  an induced unit-speed profile requires $\dot z^2+\dot r^2=1$, hence
  $|\dot r|\le1$.
\end{remark}

\begin{example}[Round and hyperbolic surfaces of revolution]
  \label{ex:pet-ch1-sphere-revolution}
  \lean{PetersenLib.sphereAsSurfaceOfRevolution}
  \source{petersen:page-0036}
  \uses{rem:pet-ch1-rotationally-symmetric, ex:pet-ch1-sphere, ex:pet-ch1-hyperbolic-space}
  \leanok
  \dcref{ch1:1.4.5}
  Revolving $R(\sin(t/R),0,\cos(t/R))$ gives $S^2(R)$ with metric
  \[
    dt^2+R^2\sin^2(t/R)d\theta^2.
  \]
  The limit $R\sin(t/R)\to t$ recovers the Euclidean polar metric. Replacing
  the trigonometric functions by their hyperbolic analogues gives
  \[
    dt^2+R^2\sinh^2(t/R)d\theta^2,
  \]
  the metric induced on $H^2(R)$ by revolving
  $R(\sinh(t/R),0,\cosh(t/R))$ in $\mathbb{R}^{2,1}$.
\end{example}

\begin{definition}[The functions $\mathrm{sn}_k$ and $\mathrm{cs}_k$]
  \label{def:pet-ch1-snk-csk}
  \lean{PetersenLib.snFunction,PetersenLib.csFunction}
  \source{petersen:page-0036,petersen:page-0037}
  \uses{rem:pet-ch1-rotationally-symmetric}
  \leanok
  Define $\mathrm{sn}_k$ and $\mathrm{cs}_k$ as the solutions of
  $\ddot x+kx=0$ with initial data
  \[
    \mathrm{sn}_k(0)=0,\quad \dot{\mathrm{sn}}_k(0)=1,
    \qquad
    \mathrm{cs}_k(0)=1,\quad \dot{\mathrm{cs}}_k(0)=0.
  \]
  Then
  \[
    \dot{\mathrm{sn}}_k=\mathrm{cs}_k,\qquad
    \dot{\mathrm{cs}}_k=-k\,\mathrm{sn}_k,\qquad
    \mathrm{cs}_k^2+k\,\mathrm{sn}_k^2=1.
  \]
  The metrics $dt^2+\mathrm{sn}_k^2(t)d\theta^2$ give
  $\mathbb{R}^2$ for $k=0$, $S^2(1/\sqrt{k})$ for $k>0$, and
  $H^2(1/\sqrt{-k})$ for $k<0$.
\end{definition}

\subsection{Polar Versus Cartesian Coordinates}

\begin{definition}[Higher-dimensional rotational symmetry]
  \label{def:pet-ch1-rot-sym-general}
  \lean{PetersenLib.rotationallySymmetricMetricHighDim}
  \source{petersen:page-0037}
  \uses{rem:pet-ch1-rotationally-symmetric}
  \leanok
  Let $ds_{n-1}^2$ be the round metric on $S^{n-1}(1)$. A rotationally
  symmetric metric on $I\times S^{n-1}$ has the form
  \[
    dt^2+\rho^2(t)ds_{n-1}^2,
  \]
  with $\rho(0)=0$ and $\rho(t)>0$ for $t>0$. Such metrics are warped products;
  their extension across $t=0$ is governed by the next criterion.
\end{definition}

\begin{lemma}[Polar-to-Cartesian metric formula]
  \label{lem:pet-ch1-polar-coordinate-change}
  \lean{PetersenLib.polarToCartesianFormula}
  \source{petersen:page-0037,petersen:page-0038}
  \uses{def:pet-ch1-rot-sym-general}
  \leanok
  Put $x=ts\in\mathbb{R}^n$, where $t=|x|>0$ and
  $s\in S^{n-1}(1)$. Then
  \[
    \sum_i(dx^i)^2=dt^2+t^2ds_{n-1}^2,
    \qquad dt=\frac1t\sum_i x^i dx^i,
  \]
  and consequently
  \[
    dt^2+\rho^2(t)ds_{n-1}^2
      =\left(\frac1{t^2}-\frac{\rho^2(t)}{t^4}\right)
         \left(\sum_i x^i dx^i\right)^2
       +\frac{\rho^2(t)}{t^2}\sum_i(dx^i)^2.
  \]
\end{lemma}
\begin{proof}
  \uses{lem:pet-ch1-polar-coordinate-change}
  \leanok
  \sketch
  Since $x^i=ts^i$,
  $dx^i=s^i dt+t\,ds^i$. The identities
  $\sum_i(s^i)^2=1$ and $\sum_i s^i ds^i=0$ give the first formula.
  Differentiating $t^2=\sum_i(x^i)^2$ gives the second. Solving for
  $ds_{n-1}^2$ and substituting proves the final identity.
\end{proof}

\begin{theorem}[Smoothness at the origin]
  \label{thm:pet-ch1-smoothness-criterion}
  \lean{PetersenLib.rotationallySymmetricSmoothnessCriterion}
  \source{petersen:page-0038,petersen:page-0039}
  \uses{lem:pet-ch1-polar-coordinate-change}
  \leanok
  For $n\ge1$, the metric
  $dt^2+\rho^2(t)ds_{n-1}^2$ extends smoothly across $t=0$ to a metric near
  the origin of $\mathbb{R}^n$ if and only if
  \[
    \rho(0)=0,
    \qquad \dot\rho(0)=1,
    \qquad \rho^{(2l)}(0)=0\quad(l\ge1).
  \]
\end{theorem}
\begin{proof}
  \leanok
  \uses{thm:pet-ch1-smoothness-criterion}
  \sketch
  By \cref{lem:pet-ch1-polar-coordinate-change}, smoothness is equivalent to
  smooth radial extensions of
  \[
    F_1=\frac{\rho^2(t)}{t^2},
    \qquad
    F_2=\frac1{t^2}-\frac{\rho^2(t)}{t^4},
  \]
  with $F_1(0)>0$. Under the displayed derivative conditions, decompose
  $\rho$ into odd and even parts. Its even part is flat at $0$. Repeated
  Hadamard division gives smooth functions $K,M,H,X$, with $K,H$ even,
  $K(0)=1$, and $M,X$ flat, such that
  \[
    \rho=t(K+M),\qquad 1-K^2=t^2H,
    \qquad M=t^2X.
  \]
  Therefore $F_1=(K+M)^2$ and $F_2=H-X(2K+M)$ are each the sum of an even
  smooth function and a flat function. Whitney's theorem for even functions
  expresses both as smooth functions of $t^2=\sum_i(x^i)^2$, giving the
  required Cartesian extensions and $F_1(0)=1$.

  Conversely, smoothness gives $F_1=1-t^2F_2$ and hence $F_1(0)=1$. On any
  line through the origin, $a(t)=F_1(te)$ is smooth and even. Near $0$ the
  odd smooth function $r(t)=t\sqrt{a(t)}$ agrees with $\rho(t)$ for $t>0$.
  Their derivatives at $0$ coincide, so $\rho(0)=0$, $\dot\rho(0)=1$, and
  all even derivatives of $\rho$ vanish.
\end{proof}

\begin{example}[Space forms]
  \label{ex:pet-ch1-space-forms}
  \lean{PetersenLib.spaceForm,PetersenLib.spaceFormIsometry,PetersenLib.spaceFormIsometry_neg}
  \source{petersen:page-0039,petersen:page-0040}
  \uses{thm:pet-ch1-smoothness-criterion, def:pet-ch1-snk-csk, ex:pet-ch1-sphere, ex:pet-ch1-hyperbolic-space}
  \leanok
  \dcref{ch1:1.4.6}
  The metrics
  $dt^2+\mathrm{sn}_k^2(t)ds_{n-1}^2$ satisfy
  \cref{thm:pet-ch1-smoothness-criterion}. Their radial intervals are
  $[0,\infty)$ for $k\le0$ and $[0,\pi/\sqrt{k}]$ for $k>0$. The resulting
  space form $S_k^n$ is isometric to $\mathbb{R}^n$ if $k=0$, to $S^n(R)$ if
  $k=1/R^2$, and to $H^n(R)$ if $k=-1/R^2$.
\end{example}
\begin{proof}
  \leanok
  \uses{ex:pet-ch1-space-forms}
  \sketch
  For $k=1/R^2$, the map
  \[
    (s,r)\longmapsto R(s\sin(r/R),\cos(r/R))
  \]
  takes $S^{n-1}\times(0,R\pi)$ into $S^n(R)$. Using
  $\sum_i(s^i)^2=1$ and $\sum_i s^i ds^i=0$, its pullback metric is
  \[
    dr^2+R^2\sin^2(r/R)ds_{n-1}^2
      =dr^2+\mathrm{sn}_{1/R^2}^2(r)ds_{n-1}^2.
  \]
  For $k=-1/R^2$, the map
  $(s,r)\mapsto(R\sinh(r/R)s,R\cosh(r/R))$ into $\mathbb{R}^{n,1}$ similarly
  pulls the Minkowski form back to
  $dr^2+R^2\sinh^2(r/R)ds_{n-1}^2$. The case $k=0$ is the Euclidean polar
  metric of \cref{ex:pet-ch1-euclidean-coords}.
\end{proof}

\subsection{Doubly Warped Products}

\begin{definition}[Doubly warped product metric]
  \label{def:pet-ch1-doubly-warped}
  \lean{PetersenLib.doublyWarpedProductMetric}
  \source{petersen:page-0041}
  \uses{def:pet-ch1-rot-sym-general}
  \leanok
  A doubly warped product metric on $I\times S^p\times S^q$ is
  \[
    dt^2+\rho^2(t)ds_p^2+\phi^2(t)ds_q^2.
  \]
  Nondegeneracy forbids simultaneous zeros of $\rho$ and $\phi$. At a zero of
  one factor, smoothness reduces to the rotational criterion for that factor,
  together with radial smoothness of the other coefficient.
\end{definition}

\begin{proposition}[Smoothness at $t=0$]
  \label{prop:pet-ch1-doubly-warped-smooth-zero}
  \lean{PetersenLib.doublyWarpedSmoothAtZero}
  \source{petersen:page-0041}
  \uses{def:pet-ch1-doubly-warped, thm:pet-ch1-smoothness-criterion}
  \leanok
  \dcref{ch1:1.4.7}
  Let $\rho,\phi:(0,b)\to(0,\infty)$ be smooth. The metric extends smoothly
  across $t=0$ if and only if
  \[
  \begin{gathered}
    \rho(0)=0,\qquad \dot\rho(0)=1,
    \qquad \rho^{(2l)}(0)=0\quad(l\ge1),\\
    \phi(0)>0,
    \qquad \phi^{(2l+1)}(0)=0\quad(l\ge0).
  \end{gathered}
  \]
  The resulting local topology is $\mathbb{R}^{p+1}\times S^q$. Positivity of
  $\phi$ near $0$ is needed in the reverse implication because the coefficient
  $\phi^2$ alone does not determine its sign.
\end{proposition}
\begin{proof}
  \leanok
  \uses{prop:pet-ch1-doubly-warped-smooth-zero,thm:pet-ch1-smoothness-criterion}
  \sketch
  In Cartesian coordinates on the collapsing $\mathbb{R}^{p+1}$ factor, the
  coefficients are
  \[
    F_1=\rho^2/t^2,\qquad F_2=1/t^2-\rho^2/t^4,
    \qquad F_3=\phi^2.
  \]
  The first two extend precisely under the conditions on $\rho$, by
  \cref{thm:pet-ch1-smoothness-criterion}. If the odd derivatives of $\phi$
  vanish, its odd part is flat and Whitney's even-function theorem makes
  $\phi(|x|)$, hence $F_3$, smooth in $x$, with positive value at the origin.
  Conversely, restrict a smooth positive extension of $F_3$ to a line.
  Its positive square root is smooth and even and agrees with $\phi$ on the
  positive half-line. Thus $\phi(0)>0$ and every odd derivative vanishes.
\end{proof}

\begin{proposition}[Smoothness at $t=b$]
  \label{prop:pet-ch1-doubly-warped-smooth-b}
  \lean{PetersenLib.doublyWarpedSmoothAtB}
  \source{petersen:page-0041}
  \uses{def:pet-ch1-doubly-warped, thm:pet-ch1-smoothness-criterion}
  \leanok
  \dcref{ch1:1.4.8}
  For smooth positive $\rho,\phi:(0,b)\to(0,\infty)$, smooth extension across
  $t=b$ is equivalent to
  \[
  \begin{gathered}
    \rho(b)=0,\qquad \dot\rho(b)=-1,
    \qquad \rho^{(2l)}(b)=0\quad(l\ge1),\\
    \phi(b)>0,
    \qquad \phi^{(2l+1)}(b)=0\quad(l\ge0).
  \end{gathered}
  \]
  The local topology is $\mathbb{R}^{p+1}\times S^q$.
\end{proposition}
\begin{proof}
  \leanok
  \uses{prop:pet-ch1-doubly-warped-smooth-b,prop:pet-ch1-doubly-warped-smooth-zero}
  \sketch
  Apply \cref{prop:pet-ch1-doubly-warped-smooth-zero} to the reflected
  coordinate $s=b-t$ and the functions $\rho(b-s),\phi(b-s)$. Reflection
  preserves even derivatives and reverses the sign of odd derivatives, turning
  the required initial slope $+1$ into $\dot\rho(b)=-1$.
\end{proof}

\begin{remark}[Topology of doubly warped completions]
  \label{rem:pet-ch1-doubly-warped-topologies}
  \lean{PetersenLib.doublyWarpedTopologyTypes}
  \source{petersen:page-0041}
  \uses{prop:pet-ch1-doubly-warped-smooth-zero, prop:pet-ch1-doubly-warped-smooth-b}
  \leanok
  If $\rho,\phi$ are positive on $(0,\infty)$ and satisfy the conditions of
  \cref{prop:pet-ch1-doubly-warped-smooth-zero}, the completed topology is
  $\mathbb{R}^{p+1}\times S^q$. On $[0,b]$, imposing the corresponding
  conditions at both endpoints produces $S^{p+1}\times S^q$. If instead
  $\rho$ collapses at one endpoint and $\phi$ at the other, the result is
  $S^{p+q+1}$.
\end{remark}

\begin{example}[Sphere as a doubly warped product]
  \label{ex:pet-ch1-sphere-doubly-warped}
  \lean{PetersenLib.sphereAsDoublyWarpedProduct}
  \source{petersen:page-0041,petersen:page-0042}
  \uses{rem:pet-ch1-doubly-warped-topologies, ex:pet-ch1-sphere}
  \leanok
  \dcref{ch1:1.4.9}
  On $[0,\pi/2]\times S^p\times S^q$, the metric
  \[
    dt^2+\sin^2(t)ds_p^2+\cos^2(t)ds_q^2
  \]
  extends to the round unit sphere $S^{p+q+1}(1)$.
\end{example}
\begin{proof}
  \leanok
  \uses{ex:pet-ch1-sphere-doubly-warped}
  \sketch
  The map
  $(t,x,y)\mapsto(x\sin t,y\cos t)$ takes the interior into the unit sphere
  of $\mathbb{R}^{p+1}\times\mathbb{R}^{q+1}$. Differentiation, using
  $x\perp T_xS^p$ and $y\perp T_yS^q$, gives the stated pullback metric; the
  endpoint criteria identify its smooth completion.
\end{proof}

\subsection{Hopf Fibrations}

The doubly warped coordinates make both the Hopf map and its higher-dimensional
quotient metrics explicit.

\begin{example}[Hopf fibration in warped coordinates]
  \label{ex:pet-ch1-hopf-revisited}
  \lean{PetersenLib.hopfFibrationRevisitedSphere,PetersenLib.pullbackForm_sphereThreeParam,PetersenLib.pullbackForm_sphereTwoParam,PetersenLib.hopfMap_sphereThreeParam,PetersenLib.hopfFibrationRevisited}
  \source{petersen:page-0042,petersen:page-0043}
  \uses{ex:pet-ch1-hopf-fibration, ex:pet-ch1-sphere-doubly-warped, ex:pet-ch1-su2-berger}
  \leanok
  \dcref{ch1:1.4.10}
  Parametrize $S^3(1)$ by
  \[
    (t,e^{i\theta_1},e^{i\theta_2})
      \longmapsto(\sin t\,e^{i\theta_1},\cos t\,e^{i\theta_2}),
  \]
  with metric
  $dt^2+\sin^2t\,d\theta_1^2+\cos^2t\,d\theta_2^2$. Parametrize
  $S^2(1/2)$ by
  \[
    (r,e^{i\theta})\longmapsto
      \big(\tfrac12\cos(2r),\tfrac12\sin(2r)e^{i\theta}\big),
  \]
  with metric $dr^2+\tfrac14\sin^2(2r)d\theta^2$. In these coordinates the
  Hopf map is
  \[
    (t,e^{i\theta_1},e^{i\theta_2})
      \longmapsto(t,e^{i(\theta_1-\theta_2)}),
  \]
  and its fibres add the same angle to $\theta_1,\theta_2$. It agrees with the
  map in \cref{ex:pet-ch1-hopf-fibration} and is a Riemannian submersion.
  Moreover,
  \[
    (t,e^{i\theta_1},e^{i\theta_2})\longmapsto
    \left[\begin{smallmatrix}
      \cos t\,e^{i\theta_1}&\sin t\,e^{i\theta_2}\\
      -\sin t\,e^{-i\theta_2}&\cos t\,e^{-i\theta_1}
    \end{smallmatrix}\right]
  \]
  is an isometry from $S^3(1)$ to $SU(2)$ with the metric of
  \cref{ex:pet-ch1-su2-berger}.
\end{example}
\begin{proof}
  \uses{ex:pet-ch1-hopf-revisited}
  \leanok
  \sketch
  The vertical vector is $\partial_{\theta_1}+\partial_{\theta_2}$. An
  orthonormal horizontal frame is
  \[
    \partial_t,\qquad
    u=\frac{\cos^2t\,\partial_{\theta_1}
              -\sin^2t\,\partial_{\theta_2}}{\cos t\sin t}.
  \]
  The target frame is
  $\partial_r,(2/\sin(2r))\partial_\theta$. The differential sends the first
  source vector to $\partial_r$ and $u$ to
  $(2/\sin(2r))\partial_\theta$, proving horizontal isometry.
\end{proof}

\begin{example}[A general doubly warped circle submersion]
  \label{ex:pet-ch1-hopf-general-submersion}
  \lean{PetersenLib.hopfFibrationGeneralSubmersion}
  \source{petersen:page-0043,petersen:page-0044}
  \uses{def:pet-ch1-riemannian-submersion, def:pet-ch1-doubly-warped}
  \leanok
  \dcref{ch1:1.4.11}
  The map
  \[
    (t,e^{i\theta_1},e^{i\theta_2})
      \longmapsto(t,e^{i(\theta_1-\theta_2)})
  \]
  is a Riemannian submersion from
  $dt^2+\rho^2d\theta_1^2+\phi^2d\theta_2^2$ to
  \[
    dr^2+\frac{(\rho\phi)^2}{\rho^2+\phi^2}d\theta^2.
  \]
\end{example}

\begin{example}[Higher-dimensional Hopf quotient]
  \label{ex:pet-ch1-hopf-higher-dim}
  \lean{PetersenLib.hopfFibrationHigherDimSphere}
  \source{petersen:page-0044}
  \uses{ex:pet-ch1-hopf-general-submersion}
  \leanok
  \dcref{ch1:1.4.12}
  Decompose the round metric on $S^{2n+1}$ as $ds_{2n+1}^2=h+g$, where $h$
  is its restriction to the Hopf distribution and $g$ to the orthogonal
  complement. Equivalently, for the orthogonal projection $\operatorname{pr}$
  onto the Hopf distribution,
  \[
  \begin{aligned}
    h(v,w)&=ds_{2n+1}^2(\operatorname{pr}v,\operatorname{pr}w),\\
    g(v,w)&=ds_{2n+1}^2(v-\operatorname{pr}v,w-\operatorname{pr}w).
  \end{aligned}
  \]
  Simultaneous scalar multiplication gives a free isometric $S^1$-action on
  $S^{2n+1}\times S^1$. The quotient map
  \[
    I\times S^{2n+1}\times S^1
      \longrightarrow I\times (S^{2n+1}\times S^1)/S^1
      \simeq I\times S^{2n+1}
  \]
  is a Riemannian submersion from
  $dt^2+\rho^2(h+g)+\phi^2d\theta^2$ to
  \[
    dt^2+\rho^2g+\frac{(\rho\phi)^2}{\rho^2+\phi^2}h.
  \]
\end{example}

\begin{example}[$S^3\simeq SU(2)$ coframe form]
  \label{ex:pet-ch1-hopf-su2-coframe}
  \lean{PetersenLib.hopfFibrationSUTwoCoframeSphere,PetersenLib.suTwoCoframe,PetersenLib.sphereMetricUnit_eq_sum_suTwoCoframe,PetersenLib.hopfComplement_eq_suTwoCoframe,PetersenLib.hopfQuotientForm_eq_suTwoCoframe}
  \source{petersen:page-0045}
  \uses{ex:pet-ch1-hopf-higher-dim, ex:pet-ch1-su2-berger}
  \leanok
  \dcref{ch1:1.4.13}
  For $n=1$, the coframe of \cref{ex:pet-ch1-su2-berger} satisfies
  \[
    ds_{S^3}^2=(\sigma^1)^2+(\sigma^2)^2+(\sigma^3)^2,
    \qquad h=(\sigma^1)^2,
    \qquad g=(\sigma^2)^2+(\sigma^3)^2.
  \]
  Hence the submersion of \cref{ex:pet-ch1-hopf-higher-dim} is
  \[
  \begin{aligned}
   &\left(I\times S^3\times S^1,
     dt^2+\rho^2\big((\sigma^1)^2+(\sigma^2)^2+(\sigma^3)^2\big)
       +\phi^2d\theta^2\right)\\
   &\quad\longrightarrow
    \left(I\times S^3,
     dt^2+\rho^2\big((\sigma^2)^2+(\sigma^3)^2\big)
       +\frac{(\rho\phi)^2}{\rho^2+\phi^2}(\sigma^1)^2\right).
  \end{aligned}
  \]
\end{example}

\begin{example}[Generalized Hopf fibration]
  \label{ex:pet-ch1-generalized-hopf}
  \lean{PetersenLib.generalizedHopfFibrationSphere,PetersenLib.generalizedHopfFibrationComplexProjective,PetersenLib.genHopfBase,PetersenLib.snocLpEquiv}
  \source{petersen:page-0045}
  \uses{ex:pet-ch1-hopf-higher-dim, ex:pet-ch1-fubini-study}
  \leanok
  \dcref{ch1:1.4.14}
  In \cref{ex:pet-ch1-hopf-higher-dim}, set
  $\rho(t)=\sin t$, $\phi(t)=\cos t$, and $0\le t\le\pi/2$. The completed
  quotient is the generalized Hopf fibration
  $S^{2n+3}\to\mathbb{CP}^{n+1}$, and the Fubini--Study metric is
  \[
    dt^2+\sin^2t\,(g+\cos^2t\,h).
  \]
\end{example}

\section{Some Tensor Concepts}

\subsection{Type Change}

The metric identifies vectors and covectors and thereby permits canonical
changes of tensor variance.

\begin{definition}[Type change]
  \label{def:pet-ch1-tensor-type-change}
  \lean{PetersenLib.tensorTypeChange}
  \source{petersen:page-0045,petersen:page-0046,petersen:page-0047}
  \uses{def:pet-ch1-riemannian-manifold}
  \leanok
  An $(s,t)$-tensor is a section of
  $TM^{\otimes s}\otimes(T^*M)^{\otimes t}$. The metric isomorphism
  \[
    TM\longrightarrow T^*M,
    \qquad v\longmapsto g(v,\cdot),
  \]
  and its inverse convert it into an $(s-k,t+k)$-tensor for every integer $k$
  such that both entries are nonnegative. For a frame $E_i$ with dual coframe
  $\sigma^i$, let
  $g_{ij}=g(E_i,E_j)$ and let $(g^{ij})=(g_{ij})^{-1}$. Type change replaces
  \[
    E_i\longleftrightarrow g_{ij}\sigma^j,
    \qquad
    \sigma^j\longleftrightarrow g^{ij}E_i.
  \]
  Thus indices are raised and lowered by $g^{ij}$ and $g_{ij}$; for an
  orthonormal frame these matrices are $\delta_{ij}$.
\end{definition}

\begin{notation}[Variance conventions for the Ricci tensor]
  \label{def:pet-ch1-ricci-types}
  \lean{PetersenLib.ricciTensorTypes}
  \source{petersen:page-0047}
  \uses{def:pet-ch1-tensor-type-change}
  \leanok
  As a $(1,1)$-tensor,
  \[
    \operatorname{Ric}(E_i)=\operatorname{Ric}^j_iE_j,
    \qquad
    \operatorname{Ric}=\operatorname{Ric}^j_iE_i\otimes\sigma^j.
  \]
  Its $(0,2)$ and $(2,0)$ forms are respectively
  \[
  \begin{aligned}
    \operatorname{Ric}
      &=\operatorname{Ric}_{jk}\sigma^j\otimes\sigma^k
        =g_{ji}\operatorname{Ric}^i_k\sigma^j\otimes\sigma^k,\\
    \operatorname{Ric}
      &=\operatorname{Ric}^{ik}E_i\otimes E_k
        =g^{jl}\operatorname{Ric}^i_jE_i\otimes E_k.
  \end{aligned}
  \]
\end{notation}

\begin{notation}[Variance conventions for the curvature tensor]
  \label{def:pet-ch1-curvature-types}
  \lean{PetersenLib.curvatureTensorTypes,PetersenLib.curvatureTwoTwoTypes}
  \source{petersen:page-0047,petersen:page-0048}
  \uses{def:pet-ch1-tensor-type-change}
  \leanok
  The $(1,3)$-curvature tensor is represented by
  \[
    R=R^i_{jkl}E_i\otimes\sigma^j\otimes\sigma^l\otimes\sigma^k.
  \]
  Lowering its vector index gives
  \[
    R=R_{ijkl}\sigma^i\otimes\sigma^j\otimes\sigma^k\otimes\sigma^l
      =R^i_{jkl}g_{il}\cdot
        \sigma^i\otimes\sigma^j\otimes\sigma^k\otimes\sigma^l,
  \]
  with the raised index $l$ conventionally placed last. Raising the last index
  gives the $(2,2)$ curvature operator
  \[
    R=R^{kl}_{ij}E_k\otimes E_l\otimes\sigma^i\otimes\sigma^j
      =R^l_{ij}g^{sk}\cdots.
  \]
  Raising another index produces a different $(2,2)$-tensor.
\end{notation}

\subsection{Contractions}

\begin{definition}[Contraction]
  \label{def:pet-ch1-contraction}
  \lean{PetersenLib.tensorContraction}
  \source{petersen:page-0048}
  \uses{def:pet-ch1-tensor-type-change}
  \leanok
  A \emph{contraction} evaluates one vector-covector pair and sums. Thus
  \[
    C(T^i_jE_i\otimes\sigma^j)=T^i_i=\operatorname{tr}T,
    \qquad C(X\otimes\omega)=\omega(X).
  \]
  A covariant tensor is first type-changed; for example,
  \[
    C(T_{ij}\sigma^i\otimes\sigma^j)=T_{ik}g^{ki}.
  \]
  Multiple contractions are defined by repeating this operation.
\end{definition}

\begin{remark}[Ricci tensor as a curvature contraction]
  \label{rem:pet-ch1-ricci-contraction}
  \lean{PetersenLib.ricciAsContraction}
  \source{petersen:page-0048}
  \uses{def:pet-ch1-contraction, def:pet-ch1-curvature-types, def:pet-ch1-ricci-types}
  \leanok
  The Ricci tensor is the contraction
  \[
    \operatorname{Ric}
      =\operatorname{Ric}^i_jE_i\otimes\sigma^j
      =R^{ij}_{ik}E_i\otimes\sigma^j
      =R^i_{iks}g^{sk}E_i\otimes\sigma^j.
  \]
  Equivalently, in covariant form,
  \[
    \operatorname{Ric}
      =\operatorname{Ric}_{ij}\sigma^i\otimes\sigma^j
      =g^{kl}R_{kijl}\sigma^i\otimes\sigma^j.
  \]
  Type change identifies these formulas.
\end{remark}

\begin{definition}[Scalar curvature]
  \label{def:pet-ch1-scalar-curvature}
  \lean{PetersenLib.scalarCurvatureOfRicci}
  \source{petersen:page-0048,petersen:page-0049}
  \uses{def:pet-ch1-contraction, def:pet-ch1-ricci-types}
  \leanok
  The \emph{scalar curvature} is the contraction of the Ricci tensor:
  \[
    \operatorname{scal}
      =\operatorname{tr}(\operatorname{Ric})
      =\operatorname{Ric}^i_i
      =R^i_{iks}g^{sk}
      =\operatorname{Ric}_{ik}g^{ki}
      =R_{ijkl}g^{ik}g^{jl}.
  \]
\end{definition}

\subsection{Inner Products of Tensors}

Two norms on linear maps will be distinguished.

\begin{definition}[Operator norm]
  \label{def:pet-ch1-operator-norm}
  \lean{PetersenLib.operatorNorm}
  \source{petersen:page-0049}
  \leanok
  For a linear map $L:V\to W$ between normed spaces,
  \[
    \|L\|=\sup_{|v|=1}|Lv|.
  \]
  If $L:V\to V$ is self-adjoint with eigenvalues
  $\lambda_1\le\cdots\le\lambda_n$, then
  $\|L\|=\max\{|\lambda_1|,|\lambda_n|\}$.
\end{definition}

\begin{definition}[Euclidean norm and trace inner product]
  \label{def:pet-ch1-euclidean-norm}
  \lean{PetersenLib.euclideanNormLinearMap,PetersenLib.traceInnerProduct}
  \source{petersen:page-0049}
  \uses{def:pet-ch1-contraction}
  \leanok
  For $L:V\to W$, define
  \[
    |L|=\sqrt{\operatorname{tr}(L^*L)}
      =\sqrt{\operatorname{tr}(LL^*)}.
  \]
  It is induced by
  \[
    \langle L_1,L_2\rangle
      =\operatorname{tr}(L_1L_2^*)
      =\operatorname{tr}(L_2L_1^*).
  \]
  For self-adjoint $L$ with eigenvalues $\lambda_i$,
  $|L|=(\sum_i\lambda_i^2)^{1/2}$. Unless specified otherwise, tensor norms
  below use this Euclidean norm rather than the operator norm.
\end{definition}

\begin{proposition}[Euclidean norm of a $(1,1)$-tensor]
  \label{prop:pet-ch1-norm-11-tensor}
  \lean{PetersenLib.euclideanNormOneOneTensor}
  \source{petersen:page-0049,petersen:page-0050}
  \uses{def:pet-ch1-euclidean-norm, def:pet-ch1-tensor-type-change}
  \leanok
  For $T=T^i_jE_i\otimes\sigma^j$, regarded as a map $TM\to TM$,
  \[
    |T|^2=T^j_iT^k_i g^{jl}g_{ki}.
  \]
  In an orthonormal frame this is $|T|^2=T^i_jT^j_i$.
\end{proposition}
\begin{proof}
  \uses{prop:pet-ch1-norm-11-tensor}
  \leanok
  \sketch
  Before type change, the dual map is
  $T^*=T^j_i\sigma^i\otimes E_j$. Raising and lowering its indices gives
  \[
    T^*=T^k_i g^{jl}g_{ki}E_j\otimes\sigma^i.
  \]
  Contracting $T\circ T^*$ as in
  \cref{def:pet-ch1-euclidean-norm} and
  \cref{def:pet-ch1-contraction} yield
  $T^i_jT^k_i g^{jl}g_{ki}$, and $g_{ij}=g^{ij}=\delta_{ij}$ gives the stated
  orthonormal expression.
\end{proof}

\begin{proposition}[Euclidean norm of a $(0,2)$-tensor]
  \label{prop:pet-ch1-norm-02-tensor}
  \lean{PetersenLib.euclideanNormZeroTwoTensor}
  \source{petersen:page-0050}
  \uses{prop:pet-ch1-norm-11-tensor}
  \leanok
  If $T=T_{ij}\sigma^i\otimes\sigma^j$, then after changing one index and
  applying \cref{prop:pet-ch1-norm-11-tensor},
  \[
    |T|^2=T_{ij}T^{ij}.
  \]
\end{proposition}

\begin{definition}[Pointwise inner product of tensors]
  \label{def:pet-ch1-pointwise-tensor-inner}
  \lean{PetersenLib.pointwiseTensorInnerProduct}
  \source{petersen:page-0050}
  \uses{def:pet-ch1-euclidean-norm}
  \leanok
  Given an orthonormal frame $E_i$ and its dual coframe $\sigma^i$, declare
  \[
    E_{i_1}\otimes\cdots\otimes E_{i_s}\otimes
    \sigma^{j_1}\otimes\cdots\otimes\sigma^{j_t}
  \]
  to be an orthonormal basis of the $(s,t)$-tensors. This determines their
  pointwise inner product and extends $g(X,Y)$ for vector fields.
\end{definition}

\begin{definition}[$L^2$ inner product of tensors]
  \label{def:pet-ch1-l2-tensor-inner}
  \lean{PetersenLib.l2TensorInnerProduct}
  \source{petersen:page-0050}
  \uses{def:pet-ch1-pointwise-tensor-inner, def:pet-ch1-volume-form-oriented}
  \leanok
  For compactly supported tensors $T_1,T_2$ of the same type, define
  \[
    (T_1,T_2)_{L^2}=\int_M g(T_1,T_2)\,\operatorname{vol}.
  \]
  This is an inner product on the space of compactly supported tensors of that
  type.
\end{definition}

\subsection{Positional Notation}

\begin{remark}[Positional index notation]
  \label{rem:pet-ch1-positional-notation}
  \lean{PetersenLib.positionalTensorNotation}
  \source{petersen:page-0050,petersen:page-0051}
  \uses{def:pet-ch1-tensor-type-change, prop:pet-ch1-norm-11-tensor}
  \leanok
  Index position records tensor-factor order. For example, a section of
  $T^*M\otimes TM\otimes T^*M$ is written
  \[
    T={T_i}^j{}_k\,\sigma^i\otimes E_j\otimes\sigma^k.
  \]
  This distinguishes $E\otimes\sigma$ from $\sigma\otimes E$. In particular,
  \[
  \begin{aligned}
    T&=T^i_jE_i\otimes\sigma^j,\\
    T^*&=T^l_k\sigma^k\otimes E_l
       =T^l_jg_{kl}g^{ij}\sigma^k\otimes E_l,
  \end{aligned}
  \]
  which recovers $|T|^2=T^i_jT^j_i$ in an orthonormal frame. Positional
  notation is not universal, so invariant notation is used whenever it
  removes ambiguity.
\end{remark}

\section{Exercises}

\begin{remark}[Cartesian product metrics]
  \label{rem:pet-ch1-ex-1}
  \lean{PetersenLib.exercise1_6_1}
  \source{petersen:page-0051,petersen:page-0052}
  \uses{def:pet-ch1-riemannian-manifold, ex:pet-ch1-flat-torus}
  \leanok
  \dcref{ch1:1.6.1}
  Equip $M\times N$ with $g_M+g_N$. Then:
  \begin{enumerate}
    \item $(\mathbb{R}^n,g_{\mathbb{R}^n})$ is the product of $n$ copies of
      $(\mathbb{R},dt^2)$;
    \item the square flat torus is
      \[
        \mathbb{R}^2/\mathbb{Z}^2
          =\big(S^1,(1/2\pi)^2d\theta^2\big)
             \times\big(S^1,(1/2\pi)^2d\theta^2\big);
      \]
    \item
      $F(\theta_1,\theta_2)=\frac1{2\pi}
        (\cos\theta_1,\sin\theta_1,\cos\theta_2,\sin\theta_2)$
      is a Riemannian embedding $T^2\to\mathbb{R}^4$.
  \end{enumerate}
\end{remark}

\begin{remark}[Quotient metric for an isometric action]
  \label{rem:pet-ch1-ex-2}
  \lean{PetersenLib.exercise1_6_2}
  \source{petersen:page-0052}
  \uses{def:pet-ch1-riemannian-submersion}
  \leanok
  \dcref{ch1:1.6.2}
  Suppose a group $G$ acts isometrically on $(M,g)$, the orbit space $M/G$ is
  a manifold, and the quotient map is a submersion. Then $M/G$ has a unique
  Riemannian metric making the quotient map a Riemannian submersion.
\end{remark}

\begin{remark}[Volume of a finite Riemannian cover]
  \label{rem:pet-ch1-ex-3}
  \lean{PetersenLib.exercise1_6_3}
  \source{petersen:page-0052}
  \uses{rem:pet-ch1-covering-metric, def:pet-ch1-volume-form-oriented}
  \leanok
  \dcref{ch1:1.6.3}
  If $M\to N$ is a Riemannian $k$-fold covering, with $k\ge1$, $N$ second
  countable, and the Riemannian measures Borel, then
  \[
    \operatorname{vol}M=k\operatorname{vol}N.
  \]
\end{remark}
\begin{proof}
  \leanok
  \uses{rem:pet-ch1-ex-3}
  \sketch
  A countable evenly covered open cover exists by second countability.
  Disjointify it into Borel sets
  $B_i=U_i\setminus\bigcup_{j<i}U_j$. Over each $B_i$, the inverse image is
  the disjoint union of $k$ sheets, and the covering is measure preserving on
  each sheet because it is a local isometry. Hence
  $\operatorname{vol}F^{-1}(B_i)=k\operatorname{vol}B_i$; countable
  additivity proves the formula. The condition $k\ge1$ excludes the vacuous
  zero-cardinality interpretation that can also admit infinite fibres.
\end{proof}

\begin{remark}[Volume form of a warped product]
  \label{rem:pet-ch1-ex-4}
  \lean{PetersenLib.exercise1_6_4}
  \source{petersen:page-0052}
  \uses{def:pet-ch1-doubly-warped, def:pet-ch1-volume-form-local}
  \leanok
  \dcref{ch1:1.6.4}
  For the metric $dr^2+\rho^2(r)g_N$ on $I\times N$, the volume form is
  \[
    \rho^{n-1}\,dr\wedge\operatorname{vol}_N.
  \]
\end{remark}

\begin{remark}[Dual frame and volume form]
  \label{rem:pet-ch1-ex-5}
  \lean{PetersenLib.exercise1_6_5}
  \source{petersen:page-0052}
  \uses{def:pet-ch1-volume-form-local}
  \leanok
  \dcref{ch1:1.6.5}
  If $E_1,\ldots,E_n$ is orthonormal, then its dual coframe is
  $\sigma^i(X)=g(E_i,X)$ and
  \[
    \operatorname{vol}=\pm\sigma^1\wedge\cdots\wedge\sigma^n,
  \]
  with sign determined by orientation.
\end{remark}

\begin{remark}[Volume form in coordinates]
  \label{rem:pet-ch1-ex-6}
  \lean{PetersenLib.exercise1_6_6}
  \source{petersen:page-0052}
  \uses{def:pet-ch1-volume-form-local, def:pet-ch1-coordinate-metric}
  \leanok
  \dcref{ch1:1.6.6}
  In coordinates $x^1,\ldots,x^n$,
  \[
    \operatorname{vol}
      =\pm\sqrt{\det[g_{ij}]}\,dx^1\wedge\cdots\wedge dx^n.
  \]
\end{remark}

\begin{remark}[Conical warped metrics]
  \label{rem:pet-ch1-ex-7}
  \lean{PetersenLib.exercise1_6_7}
  \source{petersen:page-0052}
  \uses{rem:pet-ch1-rotationally-symmetric}
  \leanok
  \dcref{ch1:1.6.7}
  Determine the geometric models of
  $dt^2+a^2t^2d\theta^2$: $a=1$ gives the Euclidean plane and $a<1$ a cone;
  describe the case $a>1$.
\end{remark}

\begin{remark}[Smoothness with a radius-$R$ sphere factor]
  \label{rem:pet-ch1-ex-8}
  \lean{PetersenLib.exercise1_6_8}
  \source{petersen:page-0052}
  \uses{thm:pet-ch1-smoothness-criterion}
  \leanok
  \dcref{ch1:1.6.8}
  For
  $dr^2+\rho^2(r)g_{S^{n-1}(R)}$ with $\rho(0)=0$, smoothness at the origin
  requires
  \[
    \dot\rho(0)=1/R,
    \qquad \rho^{(2k)}(0)=0\quad(k\ge1).
  \]
\end{remark}

\begin{remark}[Euclidean isometries as a matrix group]
  \label{rem:pet-ch1-ex-9}
  \lean{PetersenLib.exercise1_6_9}
  \source{petersen:page-0052}
  \uses{def:pet-ch1-isometry-group, ex:pet-ch1-isometry-group-euclidean}
  \leanok
  \dcref{ch1:1.6.9}
  Identify $\mathbb{R}^n$ with the hyperplane $x^{n+1}=R$ in
  $\mathbb{R}^{n+1}$. Then $\operatorname{Iso}(\mathbb{R}^n)$ is represented by
  \[
    \left[\begin{smallmatrix}O&v\\0&1\end{smallmatrix}\right],
    \qquad O\in O(n),\quad v\in\mathbb{R}^n,
  \]
  and these are precisely the linear maps preserving that hyperplane and the
  degenerate form $x^1y^1+\cdots+x^ny^n$.
\end{remark}

\begin{remark}[Sylvester basis and trace]
  \label{rem:pet-ch1-ex-10}
  \lean{PetersenLib.exercise1_6_10}
  \source{petersen:page-0052,petersen:page-0053}
  \uses{def:pet-ch1-pseudo-riemannian}
  \leanok
  \dcref{ch1:1.6.10}
  Let $g$ be a symmetric nondegenerate form of index $p$ on $V$. There is a
  basis $e_1,\ldots,e_n$ with pairwise orthogonal elements satisfying
  $g(e_i,e_i)=1$ for $i\le n-p$ and $g(e_i,e_i)=-1$ for $i>n-p$. In this basis,
  \[
    v=\sum_{i\le n-p}g(v,e_i)e_i-
      \sum_{i>n-p}g(v,e_i)e_i,
  \]
  and every linear map $L:V\to V$ satisfies
  \[
    \operatorname{tr}L
      =\sum_i\frac{g(L(e_i),e_i)}{g(e_i,e_i)}.
  \]
  Thus $(V,g)$ is linearly isometric to the corresponding $\mathbb{R}^{p,q}$.
\end{remark}

\begin{remark}[Type change by metric contraction]
  \label{rem:pet-ch1-ex-11}
  \lean{PetersenLib.exercise1_6_11}
  \source{petersen:page-0053}
  \uses{def:pet-ch1-tensor-type-change, def:pet-ch1-contraction}
  \leanok
  \dcref{ch1:1.6.11}
  Let $g^{-1}$ be the $(2,0)$ inner-product tensor on $T^*M$. Changing the type
  of a tensor is equivalent to tensoring with $g$ or $g^{-1}$ and contracting
  the appropriate pair of indices.
\end{remark}

\begin{remark}[Tensor norm from an orthonormal basis]
  \label{rem:pet-ch1-ex-12}
  \lean{PetersenLib.exercise1_6_12}
  \source{petersen:page-0053}
  \uses{def:pet-ch1-euclidean-norm}
  \leanok
  \dcref{ch1:1.6.12}
  For a $(1,1)$-tensor $T$ and an orthonormal basis $E_i$,
  \[
    |T|^2=\sum_i|T(E_i)|^2.
  \]
\end{remark}

\begin{remark}[Orthogonality of symmetric and skew tensors]
  \label{rem:pet-ch1-ex-13}
  \lean{PetersenLib.exercise1_6_13}
  \source{petersen:page-0053}
  \uses{def:pet-ch1-pointwise-tensor-inner}
  \leanok
  \dcref{ch1:1.6.13}
  If $(1,1)$-tensors $S,T$ are respectively symmetric and skew-symmetric, then
  $g(S,T)=0$.
\end{remark}

\begin{remark}[Tensor inner product by contraction]
  \label{rem:pet-ch1-ex-14}
  \lean{PetersenLib.exercise1_6_14}
  \source{petersen:page-0053}
  \uses{def:pet-ch1-pointwise-tensor-inner, def:pet-ch1-tensor-type-change, def:pet-ch1-contraction}
  \leanok
  \dcref{ch1:1.6.14}
  The inner product of two tensors of the same type is obtained by changing
  enough indices of one tensor and then contracting every remaining matched
  pair.
\end{remark}

\begin{remark}[Projective quotient away from the origin]
  \label{rem:pet-ch1-ex-15}
  \lean{PetersenLib.exercise1_6_15}
  \source{petersen:page-0053}
  \uses{def:pet-ch1-riemannian-submersion, ex:pet-ch1-fubini-study}
  \leanok
  \dcref{ch1:1.6.15}
  Let $\mathbb{F}=\mathbb{R}$ or $\mathbb{C}$ and
  $F:\mathbb{F}^{n+1}\setminus\{0\}\to\mathbb{FP}^n$ send $x$ to its
  $\mathbb{F}$-span. Give $\mathbb{FP}^n$ the metric for which the restriction
  of $F$ to the unit sphere is a Riemannian submersion. Then:
  \begin{enumerate}
    \item $F$ is a submersion, since it factors as radial projection followed
      by the Hopf quotient;
    \item $F$ is not a Riemannian submersion for the standard metric on
      $\mathbb{F}^{n+1}\setminus\{0\}$, because scale invariance makes
      horizontal vectors at $x$ and $2x$ incompatible with simultaneous
      isometry;
    \item the cylinder metric
      $dr^2+g_{S^n}$ does make $F$ a Riemannian submersion. In radial-tangential
      components it is
      \[
        g_x(u,v)=\langle u_r,v_r\rangle
          +|x|^{-2}\langle u_T,v_T\rangle.
      \]
  \end{enumerate}
  The formalization proves the obstruction in part (2) for all
  $\mathbb{F},n$ and target metrics, and proves parts (1) and (3) for
  $\mathbb{F}=\mathbb{C}$, $n=1$, namely
  $S^3\to\mathbb{CP}^1=S^2(1/2)$.
\end{remark}

\begin{remark}[Arc length]
  \label{rem:pet-ch1-ex-16}
  \lean{PetersenLib.exercise1_6_16}
  \source{petersen:page-0053,petersen:page-0054}
  \uses{def:pet-ch1-riemannian-manifold}
  \leanok
  \dcref{ch1:1.6.16}
  For a curve $c:[a,b]\to(M,g)$, define
  \[
    L(c)=\int_a^b|\dot c(t)|\,dt.
  \]
  This length is invariant under reparametrization, every curve with
  nowhere-vanishing speed admits a unit-speed parametrization, and the
  definition extends to absolutely continuous manifold-valued curves.
  The formalization includes the first two assertions; the third is omitted
  because Mathlib has no corresponding theory of absolutely continuous curves
  in manifolds.
\end{remark}

\begin{remark}[Length preservation by Riemannian immersions]
  \label{rem:pet-ch1-ex-17}
  \lean{PetersenLib.exercise1_6_17}
  \source{petersen:page-0054}
  \uses{rem:pet-ch1-ex-16, def:pet-ch1-pullback-metric}
  \leanok
  \dcref{ch1:1.6.17}
  Every Riemannian immersion preserves the length of curves.
\end{remark}

\begin{remark}[Length contraction by Riemannian submersions]
  \label{rem:pet-ch1-ex-18}
  \lean{PetersenLib.exercise1_6_18}
  \source{petersen:page-0054}
  \uses{rem:pet-ch1-ex-16, def:pet-ch1-riemannian-submersion}
  \leanok
  \dcref{ch1:1.6.18}
  If $F:(M,g_M)\to(N,g_N)$ is a Riemannian submersion and
  $c:[a,b]\to M$, then
  \[
    L(F\circ c)\le L(c).
  \]
  Equality holds exactly when
  $\dot c(t)\perp\ker DF_{c(t)}$ for all $t$.
\end{remark}

\begin{remark}[Euclidean length minimization]
  \label{rem:pet-ch1-ex-19}
  \lean{PetersenLib.exercise1_6_19}
  \source{petersen:page-0054}
  \uses{rem:pet-ch1-ex-16, ex:pet-ch1-euclidean-space}
  \leanok
  \dcref{ch1:1.6.19}
  Every curve between two points of Euclidean space has length at least the
  Euclidean distance between them. Equality is possible only when its image
  lies in the connecting segment. This follows by projecting the curve onto
  the unit vector parallel to that segment.
\end{remark}

The next algebraic obstruction supplies the flat-immersion conclusions used in
the spherical and hyperbolic exercises without importing the geodesic theory of
later chapters.

\begin{proposition}[No flat isometric immersion into an umbilic hypersurface]
  \label{prop:pet-ch1-no-flat-immersion}
  \lean{PetersenLib.no_isometricImmersion_flat_to_umbilic}
  \uses{def:pet-ch1-riemannian-manifold, def:pet-ch1-pullback-metric}
  \leanok
  Let $\beta$ be a symmetric bilinear form on a real vector space $E$ of
  dimension $n+1$, where $n\ge2$, and let $c\ne0$. No smooth map
  $G:U\to E$ from a nonempty open $U\subset\mathbb{R}^n$ can satisfy
  \[
    \beta(G,G)=c,
    \qquad
    \beta(\partial_iG,\partial_jG)=\delta_{ij}
    \quad\text{on }U.
  \]
  In particular, no open Euclidean domain admits an isometric immersion into
  either $S^n\subset\mathbb{R}^{n+1}$ or
  $H^n\subset\mathbb{R}^{n,1}$.
\end{proposition}

\begin{proof}
  \leanok
  \sketch
  Differentiating the two hypotheses gives
  \[
    \beta(\partial_iG,G)=0,
    \qquad
    \beta(\partial_i\partial_jG,G)=-\delta_{ij}.
  \]
  Set $A_{ij,k}=\beta(\partial_i\partial_jG,\partial_kG)$. Metric
  differentiation gives $A_{ij,k}=-A_{ik,j}$, while equality of mixed partials
  gives $A_{ij,k}=A_{ji,k}$. Cycling these symmetries yields
  $A_{ij,k}=-A_{ij,k}$, hence $A_{ij,k}=0$.

  The vectors $\partial_1G,\ldots,\partial_nG,G$ are $\beta$-orthogonal with
  nonzero squared lengths $1,\ldots,1,c$, so they form a basis of $E$.
  The preceding pairings therefore determine
  \[
    \partial_i\partial_jG=-\frac{\delta_{ij}}cG.
  \]
  For distinct indices $1,2$, put $\Phi=\partial_2G$. Symmetry of
  $\partial_w\partial_v\Phi$ in $v,w$, evaluated at $v=e_2,w=e_1$, forces
  $-c^{-1}\partial_1G=0$, contradicting
  $\beta(\partial_1G,\partial_1G)=1$.
\end{proof}

\begin{remark}[Great-circle length minimization]
  \label{rem:pet-ch1-ex-20}
  \lean{PetersenLib.exercise1_6_20}
  \source{petersen:page-0054}
  \uses{ex:pet-ch1-sphere, rem:pet-ch1-ex-16, prop:pet-ch1-no-flat-immersion}
  \leanok
  \dcref{ch1:1.6.20}
  Fix $p,q\in S^n\subset\mathbb{R}^{n+1}$ and a unit
  $v\in T_pS^n$. Then:
  \begin{enumerate}
    \item $p\cos t+v\sin t$ is unit speed with initial data $(p,v)$;
    \item
      $F(r,v)=p\cos r+v\sin r$ is a diffeomorphism
      $(0,\pi)\times S^{n-1}\to S^n\setminus\{\pm p\}$;
    \item if $q=F(r_0,v_0)$, its radial field satisfies
      \[
        \partial_r|_q
          =\frac{-p+(p\cdot q)q}{\sqrt{1-(p\cdot q)^2}}
          =-p\sin r_0+v_0\cos r_0;
      \]
    \item integrating $\dot c\cdot\partial_r=dr/dt$ shows that a curve from
      $p$ to $q$ has length at least $r_0$, with equality only along the
      great circle;
    \item no open $U\subset\mathbb{R}^n$ admits a Riemannian immersion into
      $S^n$; the source derives this from the fact that equilateral spherical
      triangles have angles greater than $\pi/3$.
  \end{enumerate}
  The formalization proves the last assertion through
  \cref{prop:pet-ch1-no-flat-immersion}, avoiding the later result that local
  isometries map line segments to target geodesics.
\end{remark}

\begin{remark}[Hyperbolic length minimization]
  \label{rem:pet-ch1-ex-21}
  \lean{PetersenLib.exercise1_6_21}
  \source{petersen:page-0054,petersen:page-0055}
  \uses{ex:pet-ch1-hyperbolic-space, rem:pet-ch1-ex-16, prop:pet-ch1-no-flat-immersion}
  \leanok
  \dcref{ch1:1.6.21}
  Fix $p,q\in H^n\subset\mathbb{R}^{n,1}$ and a unit
  $v\in T_pH^n$, so $g(p,p)=g(q,q)=-1$, $g(v,v)=1$, and $g(p,v)=0$. Then:
  \begin{enumerate}
    \item $p\cosh t+v\sinh t$ is unit speed with initial data $(p,v)$;
    \item
      $F(r,v)=p\cosh r+v\sinh r$ is a diffeomorphism
      $(0,\infty)\times S^{n-1}\to H^n\setminus\{p\}$;
    \item if $q=F(r_0,v_0)$, then
      \[
        \partial_r|_q
          =\frac{-p-(q\cdot p)q}{\sqrt{-1+(q\cdot p)^2}}
          =p\sinh r_0+v_0\cosh r_0;
      \]
    \item every curve from $p$ to $q$ has length at least $r_0$, with equality
      only along the corresponding hyperbola;
    \item no open $U\subset\mathbb{R}^n$ admits a Riemannian immersion into
      $H^n$; the source uses equilateral hyperbolic triangles, whose angles are
      less than $\pi/3$.
  \end{enumerate}
  The formalization obtains the final assertion from
  \cref{prop:pet-ch1-no-flat-immersion} with $\beta$ equal to the Minkowski
  form and $c=-1$.
\end{remark}

\begin{remark}[Quaternionic Hopf fibrations]
  \label{rem:pet-ch1-ex-22}
  \lean{PetersenLib.exercise1_6_22}
  \source{petersen:page-0055,petersen:page-0056}
  \uses{ex:pet-ch1-hopf-fibration}
  \leanok
  \dcref{ch1:1.6.22}
  Write $q=a+bi+cj+dk=z+wj$ and identify
  \[
    q\longleftrightarrow
      \left[\begin{smallmatrix}z&w\\-\bar w&\bar z\end{smallmatrix}\right].
  \]
  This gives the associative real-bilinear quaternion product. With
  $\bar q=a-bi-cj-dk$,
  \[
    |q|^2=q\bar q=\bar qq=|z|^2+|w|^2
      =\det\left[\begin{smallmatrix}z&w\\-\bar w&\bar z\end{smallmatrix}\right],
  \]
  $|pq|=|p||q|$, and $\overline{pq}=\bar q\bar p$. Define
  \[
  \begin{aligned}
    H^l(p,q)&=\big(\tfrac12(|p|^2-|q|^2),\bar p q\big),\\
    H^r(p,q)&=\big(\tfrac12(|p|^2-|q|^2),p\bar q\big).
  \end{aligned}
  \]
  Both maps send $S^7(1)\subset\mathbb{H}^2$ to
  $S^4(1/2)\subset\mathbb{R}\oplus\mathbb{H}$. The fibres of $H^l$ and
  $H^r$ are respectively the left and right diagonal orbits of unit
  quaternions, and both maps are Riemannian submersions.
\end{remark}

\begin{remark}[Euler-number bundles over $S^2$]
  \label{rem:pet-ch1-ex-23}
  \lean{PetersenLib.exercise1_6_23}
  \source{petersen:page-0056,petersen:page-0057}
  \uses{ex:pet-ch1-hopf-su2-coframe, thm:pet-ch1-smoothness-criterion}
  \dcref{ch1:1.6.23}
  In the submersion of \cref{ex:pet-ch1-hopf-su2-coframe}, assume
  $\rho,\phi>0$ on $(0,\infty)$ and set
  \[
    f=\rho,
    \qquad
    h=\frac{(\rho\phi)^2}{\rho^2+\phi^2}.
  \]
  If
  \[
    f(0)>0,
    \quad f^{(2l+1)}(0)=0\ (l\ge0),
    \quad h(0)=0,
    \quad h'(0)=k,
    \quad h^{(2l)}(0)=0\ (l\ge1)
  \]
  for a positive integer $k$, the quotient metric extends smoothly to the
  vector bundle over $S^2$ with Euler number $\pm k$. Away from the zero
  section it is $(0,\infty)\times S^3/\mathbb{Z}_k$, obtained as a quotient of
  $S^3\times\mathbb{R}^2$. For $k=2$ this is $TS^2$; the source describes the
  $k=1$ case as looking like $\mathbb{CP}^2\setminus\{\mathrm{pt}\}$.
  The formalization covers positivity,
  the Riemannian submersion, and analytic closing at the origin; it does not
  identify the bundle topology, Euler number, or these two special cases.
\end{remark}

\begin{remark}[Bi-invariant metrics on compact Lie groups]
  \label{rem:pet-ch1-ex-24}
  \lean{PetersenLib.exercise1_6_24}
  \source{petersen:page-0057}
  \uses{def:pet-ch1-left-invariant-metric, def:pet-ch1-volume-form-local}
  \leanok
  \dcref{ch1:1.6.24}
  Let $G$ be compact and $g_l$ left-invariant. Averaging over right
  translations gives the bi-invariant metric
  \[
    g(v,w)=\left(\int_G\operatorname{vol}\right)^{-1}
      \int_G g_l(DR_xv,DR_xw)\,\operatorname{vol}.
  \]
  For any bi-invariant metric, conjugation by $h\in G$ is an isometry, so
  $\operatorname{Ad}_h:\mathfrak g\to\mathfrak g$ is a linear isometry.
  Differentiating this invariance gives
  \[
    g([U,X],Y)=-g(X,[U,Y]),
  \]
  so $\operatorname{ad}_U$ is skew-symmetric.
\end{remark}

\begin{remark}[Characterization of bi-invariant pseudo-metrics]
  \label{rem:pet-ch1-ex-25}
  \lean{PetersenLib.exercise1_6_25}
  \source{petersen:page-0057}
  \uses{def:pet-ch1-pseudo-riemannian, def:pet-ch1-left-invariant-metric}
  \leanok
  \dcref{ch1:1.6.25}
  A nondegenerate symmetric bilinear form $(\cdot,\cdot)$ on the Lie algebra
  $\mathfrak g$ defines a bi-invariant pseudo-Riemannian metric on $G$ if and
  only if
  \[
    (X,Y)=(\operatorname{Ad}_hX,\operatorname{Ad}_hY)
    \quad\text{for every }h\in G.
  \]
\end{remark}

\begin{remark}[Averaging to an invariant metric]
  \label{rem:pet-ch1-ex-26}
  \lean{PetersenLib.exercise1_6_26}
  \source{petersen:page-0058}
  \uses{def:pet-ch1-riemannian-manifold}
  \leanok
  \dcref{ch1:1.6.26}
  If a compact Lie group $G$ acts smoothly on $M$, then $M$ admits a metric
  for which the action is isometric. Start with a metric constructed by a
  partition of unity and average it against normalized Haar measure.
  In local coordinates, all derivatives of the averaged coefficients may be
  passed under the integral because the differentiated integrands are jointly
  continuous on the compact group; hence the averaged metric is smooth.
\end{remark}

\begin{remark}[Killing form]
  \label{rem:pet-ch1-ex-27}
  \lean{PetersenLib.exercise1_6_27}
  \source{petersen:page-0058}
  \uses{def:pet-ch1-left-invariant-metric}
  \leanok
  \dcref{ch1:1.6.27}
  The Killing form on $\mathfrak g$ is
  \[
    B(X,Y)=\operatorname{tr}(\operatorname{ad}_X\circ\operatorname{ad}_Y).
  \]
  It is symmetric and bilinear. If $G$ admits a bi-invariant metric, then
  $B(X,X)\le0$, and
  \[
    B(\operatorname{ad}_ZX,Y)=-B(X,\operatorname{ad}_ZY).
  \]
  For connected $G$, the function
  $t\mapsto B(\operatorname{Ad}_{\exp(tZ)}X,
              \operatorname{Ad}_{\exp(tZ)}Y)$ is constant, so $B$ is
  $\operatorname{Ad}$-invariant. When $\mathfrak g$ is semisimple, $B$ is
  nondegenerate and defines a bi-invariant pseudo-Riemannian metric; if $G$ is
  also compact, $-B$ is Riemannian.
\end{remark}

\begin{remark}[Bi-invariant pseudo-metric on $\mathrm{SL}(n,\mathbb{R})$]
  \label{rem:pet-ch1-ex-28}
  \lean{PetersenLib.exercise1_6_28}
  \source{petersen:page-0058}
  \uses{def:pet-ch1-pseudo-riemannian, def:pet-ch1-left-invariant-metric}
  \leanok
  \dcref{ch1:1.6.28}
  On the trace-zero Lie algebra $\mathfrak{sl}(n,\mathbb{R})$, the form
  \[
    (X,Y)=\operatorname{tr}(XY)
  \]
  is symmetric, nondegenerate, and $\operatorname{Ad}$-invariant. It therefore
  defines a bi-invariant pseudo-Riemannian metric on
  $\operatorname{SL}(n,\mathbb{R})$.
\end{remark}

\begin{remark}[A group without a bi-invariant pseudo-metric]
  \label{rem:pet-ch1-ex-29}
  \lean{PetersenLib.exercise1_6_29}
  \source{petersen:page-0058}
  \uses{def:pet-ch1-pseudo-riemannian, def:pet-ch1-left-invariant-metric}
  \leanok
  \dcref{ch1:1.6.29}
  The matrices of $\operatorname{diag}(a^{-1},ab,1)$-type
  \[
    \left[\begin{smallmatrix}
      a^{-1}&0&0\\0&ab&\ast\\0&0&1
    \end{smallmatrix}\right],
    \qquad a>0,\quad b\in\mathbb{R},
  \]
  form a two-dimensional Lie group that admits no bi-invariant
  pseudo-Riemannian metric.
\end{remark}
